% Unit 076 English translation of chapter6.tex lines 869--1093.
% Source: Methods of Algebra, Volume 2, commit
% 9a5803ff77dd3257484cb177f851a73770a59dd3 (CC BY 4.0).
% stable-unit-id: o014.aljabr2.chapter6.change-of-group
% Coverage: Section 6.5 in full; ends before Section 6.6 at line 1094.

% segment-id: o014.li2.chapter6.change-of-group.h001
\section{Change of groups}\label{sec:change-of-group}
\hypertarget{o014.aljabr2.chapter6.change-of-group}{ }

% segment-id: o014.li2.chapter6.change-of-group.p001
Consider a group homomorphism $\varphi:H\to G$. In what follows, $M$ denotes
an arbitrary $G$-module and $N$ an arbitrary $H$-module. Our first task in
this section is to define two families of canonical homomorphisms
% segment-id: o014.li2.chapter6.change-of-group.e001
\begin{equation}\label{eqn:change-of-group-coh}\begin{aligned}
	\Hm^n(G, M) & \to \Hm^n(H, \varphi^* M), \\
	\Hm_n(H, \varphi^* M) & \to \Hm_n(G, M),
\end{aligned}\end{equation}
where $n\in\Z_{\geq0}$. Since
$\varphi^*:G\dcate{Mod}\to H\dcate{Mod}$ is exact, the functors on both
sides of these homomorphisms have long exact sequences; in other words, they
are cohomological or homological $\delta$-functors. The homomorphisms above
will be compatible with the long exact sequences. We present two
constructions.

% segment-id: o014.li2.chapter6.change-of-group.p002
For cohomology, the first method considers the family of functors
$B^n:=\Hm^n(H,\varphi^*(\cdot))$. There is a canonical homomorphism
% segment-id: o014.li2.chapter6.change-of-group.q001
\[ M^G \hookrightarrow M^{\varphi(H)} = (\varphi^* M)^H =  B^0(M). \]
Because $(\Hm^n(G,\cdot))_n$ is a universal cohomological $\delta$-functor,
this homomorphism extends uniquely to a morphism of cohomological
$\delta$-functors $\Hm^n(G,\cdot)\to B^n$.

% segment-id: o014.li2.chapter6.change-of-group.p003
The homology case is similar; the key is to use the canonical homomorphism
$(\varphi^*M)_H=M_{\varphi(H)}\twoheadrightarrow M_G$.

% segment-id: o014.li2.chapter6.change-of-group.p004
The second method works in the derived category. We introduce some related
notation.

% segment-id: o014.li2.chapter6.change-of-group.c001
\begin{convention}
	Write $\cate{C}(G):=\cate{C}(G\dcate{Mod})$,
	$\cate{K}(G):=\cate{K}(G\dcate{Mod})$, and
	$\cate{D}(G):=\cate{D}(G\dcate{Mod})$; use similar notation when
	one-sided or two-sided boundedness conditions are imposed. Likewise, write
	$\cate{C}(\Bbbk):=\cate{C}(\Bbbk\dcate{Mod})$, together with
	$\cate{K}(\Bbbk)$, $\cate{D}(\Bbbk)$, and so forth.
\end{convention}

% segment-id: o014.li2.chapter6.change-of-group.p005
The functor on derived categories induced by the exact functor $\varphi^*$ is
still denoted
$\varphi^*:\cate{D}(G)\to\cate{D}(H)$. Use the customary notation
$\mathrm{R}$ and $\mathrm{L}$ for right- and left-derived functors,
respectively. We seek morphisms between triangulated functors
% segment-id: o014.li2.chapter6.change-of-group.q002
\begin{align*}
	\mathrm{R}\left( (\cdot)^G \right) & \to \mathrm{R}\left( (\cdot)^H \right) \circ \varphi^* , \\
	\mathrm{L}\left( (\cdot)_H \right) \circ \varphi^* & \to \mathrm{L}\left((\cdot)_G \right).
\end{align*}

% segment-id: o014.li2.chapter6.change-of-group.p006
The key is the universal property. In the cohomology case, lift the functors
$(\cdot)^G$, $(\cdot)^H$, and so on to the level of complexes, with the same
notation. Consider the following 2-cell diagram involving categories,
functors, and morphisms of functors, all understood to be triangulated.
% segment-id: o014.li2.chapter6.change-of-group.g001
\begin{equation}\label{eqn:grp-pullback-2-cell}
	\begin{tikzcd}[column sep=huge, row sep=large]
		\cate{K}^+(G) \arrow[d] \arrow[r, "{\varphi^*}"] \arrow[rr, bend left, "{(\cdot)^G}", "{}"' {name=U}] & \cate{K}^+(H) \arrow[d] \arrow[r, "{(\cdot)^H}"] \arrow[from=U, Rightarrow] \arrow[Rightarrow, ld, "\sim" sloped] & \cate{K}^+(\Bbbk) \arrow[d] \arrow[Rightarrow, ld] \\
		\cate{D}^+(G) \arrow[r, "{\varphi^*}"'] & \cate{D}^+(H) \arrow[r, "{\mathrm{R}((\cdot)^H)}"'] & \cate{D}^+(\Bbbk)
	\end{tikzcd}
\end{equation}
The $\Rightarrow$ on the curved part comes from the previously used morphism
$(\cdot)^G\to(\cdot)^H\circ\varphi^*$. The
$\stackrel{\sim}{\Rightarrow}$ on the left square comes from the exactness of
$\varphi^*$, while the $\Rightarrow$ on the right square is part of the data
of the right-derived functor $\mathrm{R}((\cdot)^H)$. The superscript $+$ is
used here only for simplicity; the unbounded version requires the theory of
\S\ref{sec:unbounded-derived}.

% segment-id: o014.li2.chapter6.change-of-group.p007
Recall that 2-cell diagrams can be composed vertically, meaning that the
morphisms $\Rightarrow$ between functors are composed. By the universal
property of right-derived functors (see
Remark~\ref{rem:derived-can-morphism} and the preceding discussion), there is
a unique $\Downarrow$ in the left-hand diagram below such that
% segment-id: o014.li2.chapter6.change-of-group.g002
	{\scriptsize
	\begin{equation}\label{eqn:change-of-groups-univ}
	\begin{tikzcd}[column sep=huge, row sep=large]
		\cate{K}^+(G) \arrow[r, "{(\cdot)^G}"] \arrow[d] & \cate{K}^+(\Bbbk) \arrow[d] \arrow[Rightarrow, ld] \\
		\cate{D}^+(G) \arrow[r, "{\mathrm{R}((\cdot)^G)}", ""' {name=D}] \arrow[r, bend right=60, "" {name=U}, "{\mathrm{R}((\cdot)^H) \circ \varphi^*}"'] & \cate{D}^+(\Bbbk) \arrow[Rightarrow, from=D, to=U, "{\exists!}"']
	\end{tikzcd} \xlongequal{\text{vertical composite}} \;\text{vertical composite of \eqref{eqn:grp-pullback-2-cell}}
	\begin{tikzcd}[column sep=huge, row sep=large]
		\cate{K}^+(G) \arrow[r, "{(\cdot)^G}"] \arrow[d] & \cate{K}^+(\Bbbk) \arrow[d] \arrow[Rightarrow, ld] \\
		\cate{D}^+(G) \arrow[r, "{\mathrm{R}((\cdot)^H) \circ \varphi^*}"' inner sep=0.6em] & \cate{D}^+(\Bbbk)
	\end{tikzcd}
	\end{equation}
	}
The symbol \rotatebox[origin=c]{45}{$\Leftarrow$} in the left-hand diagram is
part of the data of the right-derived functor
$\mathrm{R}\left((\cdot)^G\right)$. This uniquely determines the morphism
% segment-id: o014.li2.chapter6.change-of-group.q003
\[ \mathrm{R}\left( (\cdot)^G \right) \to \mathrm{R}\left((\cdot)^H \right) \circ \varphi^* . \]

% segment-id: o014.li2.chapter6.change-of-group.p008
The method for homology is dual. It uses the universal property of
left-derived functors and the previously used morphism
$(\cdot)_H\circ\varphi^*\to(\cdot)_G$. We next give a more direct
description.

% segment-id: o014.li2.chapter6.change-of-group.le001
\begin{lemma}\label{prop:change-of-groups-alt}
	For every group homomorphism $\varphi:H\to G$, the morphisms characterized
	above can be defined as follows.
% segment-id: o014.li2.chapter6.change-of-group.l001
	\begin{itemize}
% segment-id: o014.li2.chapter6.change-of-group.i001
		\item In the cohomology case, take the composite
% segment-id: o014.li2.chapter6.change-of-group.e002
		\begin{equation}\label{eqn:change-of-group-dec-1}
			\mathrm{R}\left( (\cdot)^G \right) \to \mathrm{R}\left( (\cdot)^H \circ \varphi^* \right) \to \mathrm{R}\left( (\cdot)^H \right) \circ \varphi^* .
		\end{equation}
		The first part comes from
		$(\cdot)^G\to(\cdot)^H\circ\varphi^*$, since a morphism of functors
		induces a morphism between their right-derived functors; the second
		part comes from
		Theorem~\ref{prop:localization-triangulated-composite}~(i).
% segment-id: o014.li2.chapter6.change-of-group.i002
		\item The homology case is similar: take the composite
% segment-id: o014.li2.chapter6.change-of-group.e003
		\begin{equation}\label{eqn:change-of-group-dec-2}
			\mathrm{L}\left( (\cdot)_H \right) \circ \varphi^* \to \mathrm{L} \left( (\cdot)_H \circ \varphi^* \right) \to \mathrm{L}\left( (\cdot)_G \right).
		\end{equation}
		The first part comes from
		Theorem~\ref{prop:localization-triangulated-composite}~(ii), and the
		second from $(\cdot)_H\circ\varphi^*\to(\cdot)_G$.
	\end{itemize}
\end{lemma}
% segment-id: o014.li2.chapter6.change-of-group.pf001
\begin{proof}
	It suffices to discuss the cohomology case. First, for the morphisms in
	\eqref{eqn:change-of-group-dec-1}, check the following equalities of
	vertical composites:
% segment-id: o014.li2.chapter6.change-of-group.g003
	\[\begin{tikzcd}[column sep=huge, row sep=large]
		\cate{K}^+(G) \arrow[r, "{(\cdot)^G}"] \arrow[d] & \cate{K}^+(\Bbbk) \arrow[d] \arrow[Rightarrow, ld] \\
		\cate{D}^+(G) \arrow[r, "{\mathrm{R}((\cdot)^G)}", ""' {name=D}] \arrow[r, bend right=60, "" {name=U}, "{\mathrm{R}((\cdot)^H \circ \varphi^*)}"'] & \cate{D}^+(\Bbbk) \arrow[Rightarrow, from=D, to=U]
	\end{tikzcd} = \begin{tikzcd}[column sep=huge, row sep=large]
		\cate{K}^+(G) \arrow[r, "{(\cdot)^H \circ \varphi^*}" {name=D}] \arrow[d] \arrow[r, bend left=60, ""' {name=U}, "{(\cdot)^G}"] & \cate{K}^+(\Bbbk) \arrow[d] \arrow[Rightarrow, ld] \\
		\cate{D}^+(G) \arrow[r, "{\mathrm{R}((\cdot)^H \circ \varphi^* ) }"'] & \cate{D}^+(\Bbbk) \arrow[Rightarrow, from=U, to=D]
	\end{tikzcd}\]
and
% segment-id: o014.li2.chapter6.change-of-group.g004
	\[\begin{tikzcd}[column sep=huge, row sep=large]
		\cate{K}^+(G) \arrow[r, "{(\cdot)^H \circ \varphi^*}"] \arrow[d] & \cate{K}^+(\Bbbk) \arrow[d] \arrow[Rightarrow, ld] \\
		\cate{D}^+(G) \arrow[r, "{\mathrm{R}((\cdot)^H \circ \varphi^* ) }"' {name=U}] \arrow[r, bend right=60, "{\mathrm{R}((\cdot)^H) \circ \varphi^*}"', "" {name=D}] & \cate{D}^+(\Bbbk) \arrow[Rightarrow, from=U, to=D]
	\end{tikzcd} = \begin{tikzcd}[column sep=huge, row sep=large]
		\cate{K}^+(G) \arrow[d] \arrow[r, "{\varphi^*}"] & \cate{K}^+(H) \arrow[d] \arrow[r, "{(\cdot)^H}"] \arrow[Rightarrow, ld, "\sim" sloped] & \cate{K}^+(\Bbbk) \arrow[d] \arrow[Rightarrow, ld] \\
		\cate{D}^+(G) \arrow[r, "{\varphi^*}"'] & \cate{D}^+(H) \arrow[r, "{\mathrm{R}((\cdot)^H)}"'] & \cate{D}^+(\Bbbk).
	\end{tikzcd}\]

	The first equality characterizes the induced morphism
	$\mathrm{R}((\cdot)^G)\to
	\mathrm{R}((\cdot)^H\circ\varphi^*)$, while the second characterizes the
	canonical morphism in
	Theorem~\ref{prop:localization-triangulated-composite}. Both arise from
	the universal property of right-derived functors; see the explanation at
	the beginning of \S\ref{sec:triangulated-functor-localization}. Splicing
	these equalities together shows that the composite
	\eqref{eqn:change-of-group-dec-1} satisfies
	\eqref{eqn:change-of-groups-univ}.
\end{proof}

% segment-id: o014.li2.chapter6.change-of-group.r001
\begin{remark}\label{rem:change-of-group-Yoneda}
	The first method, based on the universal cohomological
	$\delta$-functor property, can be combined with derived-category theory
	to give another concise construction for the cohomology case of
	\eqref{eqn:change-of-group-coh}. Recall that for every $G$-module $M$ and
	every $n$, there are canonical isomorphisms
% segment-id: o014.li2.chapter6.change-of-group.q004
	\[ \Hm^n(G, M) \simeq \Ext^n_G(\Bbbk, M) \simeq \Hom_{\cate{D}(G)}(\Bbbk, M[n]), \quad \Ext^\bullet_G := \Ext^\bullet_{\Bbbk[G]}; \]
	and likewise for $H$-modules. Since
	$\varphi^*\Bbbk\simeq\Bbbk$ and $\varphi^*$ is a triangulated functor, we
	assert that the desired canonical homomorphism is the composite
% segment-id: o014.li2.chapter6.change-of-group.q005
	\[ \Hom_{\cate{D}(G)}(\Bbbk, M[n]) \xrightarrow{\text{functor} \;\varphi^*} \Hom_{\cate{D}(H)}\left( \varphi^* \Bbbk, \varphi^* (M[n])\right) \simeq \Hom_{\cate{D}(H)}\left(\Bbbk, (\varphi^* M)[n]\right). \]

	To prove this, it suffices to note that when $n=0$ the formula becomes the
	evident embedding $M^G\hookrightarrow(\varphi^*M)^H$ (because
	$\varphi^*\Bbbk\simeq\Bbbk$ sends $1$ to $1$), and that every morphism in
	the composite above is compatible with the long exact sequence in the
	second variable of the $\Hom$ functor in a triangulated category
	(Proposition~\ref{prop:cohomological-Hom-functor}).
\end{remark}

% segment-id: o014.li2.chapter6.change-of-group.r002
\begin{remark}\label{rem:change-of-group-coh}
	If $\varphi$ is the inclusion of a subgroup, then
	$\varphi^*=\Res^G_H$ preserves injective and projective objects. By
	Theorem~\ref{prop:localization-triangulated-composite}~(iii) and (iv), the
	two morphisms in Lemma~\ref{prop:change-of-groups-alt},
% segment-id: o014.li2.chapter6.change-of-group.q006
	\[ \mathrm{R}\left( (\cdot)^H \circ \varphi^* \right) \to \mathrm{R}\left( (\cdot)^H \right) \circ \varphi^*, \quad
	\mathrm{L}\left( (\cdot)_H \right) \circ \varphi^* \to \mathrm{L} \left( (\cdot)_H \circ \varphi^* \right) \]
	are isomorphisms. The description of the canonical morphisms therefore
	simplifies further: take an injective resolution
	$M\to I:=[I^0\to I^1\to\cdots]$. Then
	$\Res^G_H(M)\to\Res^G_H(I)$ is also an injective resolution, and the
	action on $M$ of the morphism
	$\mathrm{R}\left((\cdot)^G\right)\to
	\mathrm{R}\left((\cdot)^H\right)\circ\Res^G_H$ is
% segment-id: o014.li2.chapter6.change-of-group.q007
	\[ I^G \to \left( \Res^G_H(I)\right)^H, \quad I^G := \left[ I^{0, G} \to I^{1, G} \to \cdots \right], \;\text{and similarly}. \]

	In the homology case, take a projective resolution $P\to M$. The action on
	$M$ of the morphism
	$\mathrm{L}\left((\cdot)_H\right)\circ\Res^G_H
	\to\mathrm{L}\left((\cdot)_G\right)$ is
	$\left(\Res^G_H(P)\right)_H\to P_G$.
\end{remark}

% segment-id: o014.li2.chapter6.change-of-group.le002
\begin{lemma}\label{prop:change-of-group-compo}
	Let $F\xleftarrow{\psi}G\xleftarrow{\varphi}H$ be group homomorphisms,
	$L$ an $F$-module, and $n\in\Z_{\geq0}$. Then the composite
	$\Hm^n(F,L)\to\Hm^n(G,\psi^*L)\to
	\Hm^n(H,\varphi^*\psi^*L)$ equals
	$\Hm^n(F,L)\to\Hm^n(H,(\psi\varphi)^*L)$.

	Likewise, the composite
	$\Hm_n(H,\varphi^*\psi^*L)\to\Hm_n(G,\psi^*L)\to\Hm_n(F,L)$ equals
	$\Hm_n(H,(\psi\varphi)^*L)\to\Hm_n(F,L)$.

	Analogous equalities also hold at the level of derived categories.
\end{lemma}
% segment-id: o014.li2.chapter6.change-of-group.pf002
\begin{proof}
	The assertion about $\Hm^n$ (or $\Hm_n$) is an immediate application of
	the universal cohomological (or homological) $\delta$-functor property
	and need only be checked for $n=0$.

	For the derived-category version, the assertion about $\Hm^n$ is replaced
	by
% segment-id: o014.li2.chapter6.change-of-group.q008
	\[ \left[ \mathrm{R}((\cdot)^F) \to \mathrm{R}((\cdot)^G) \circ \psi^* \to \mathrm{R}((\cdot)^H) \circ \varphi^* \circ \psi^* \right] \xlongequal{\text{composite}} \left[ \mathrm{R}((\cdot)^F) \to \mathrm{R}((\cdot)^H) \circ (\psi\varphi)^* \right]. \]
	For this, it suffices to check the property
	\eqref{eqn:change-of-groups-univ} for the left-hand side. This reduces to
	an interesting and straightforward exercise in the composition of
	2-cells, which the reader is invited to try.
\end{proof}

% segment-id: o014.li2.chapter6.change-of-group.p009
We now generalize the canonical homomorphisms
\eqref{eqn:change-of-group-coh} further.

% segment-id: o014.li2.chapter6.change-of-group.d001
\begin{definition}\label{def:group-equivariance}
	Let $M$ be a $G$-module and $N$ an $H$-module. Below, $f$ is always a
	$\Bbbk$-module homomorphism and $h$ is an arbitrary element of $H$.
% segment-id: o014.li2.chapter6.change-of-group.l002
	\begin{itemize}
% segment-id: o014.li2.chapter6.change-of-group.i003
		\item If $f:M\to N$ satisfies $f(\varphi(h)x)=hf(x)$ for $x\in M$,
		then $f$ is called \emph{$\varphi$-equivariant}.
% segment-id: o014.li2.chapter6.change-of-group.i004
		\item If $f:N\to M$ satisfies $f(hy)=\varphi(h)f(y)$ for $y\in N$,
		then $f$ is called \emph{$\varphi$-equivariant}.
	\end{itemize}
\end{definition}

% segment-id: o014.li2.chapter6.change-of-group.p010
Equivariance is also equivalent to saying that $f$ gives an $H$-module
homomorphism $\varphi^*M\to N$ or $N\to\varphi^*M$. Moreover,
$\identity_M:M\to\varphi^*M$ and
$\identity_M:\varphi^*M\to M$ are both $\varphi$-equivariant.

% segment-id: o014.li2.chapter6.change-of-group.d002
\begin{definition}\label{def:change-of-group-equiv}
% segment-id: o014.li2.chapter6.change-of-group.x001
	\index{equivariant homomorphism}
	If $f:M\to N$ is $\varphi$-equivariant, define a family of canonical
	homomorphisms
% segment-id: o014.li2.chapter6.change-of-group.q009
	\[ \Hm^n(f) := \left[ \Hm^n(G, M) \to \Hm^n(H, \varphi^* M) \to \Hm^n(H, N) \right] \;\text{as the composite}; \]
	if $f:N\to M$ is $\varphi$-equivariant, define a family of canonical
	homomorphisms
% segment-id: o014.li2.chapter6.change-of-group.q010
	\[ \Hm_n(f) := \left[ \Hm_n(H, N) \to \Hm_n(H, \varphi^* M) \to \Hm_n(G, M) \right] \;\text{as the composite}. \]
\end{definition}

% segment-id: o014.li2.chapter6.change-of-group.p011
Notice that $f$ and $\varphi$ point in opposite directions in the cohomology
case and in the same direction in the homology case. The homomorphisms in
Definition~\ref{def:change-of-group-equiv} are compatible with long exact
sequences. They include the functoriality of $\Hm^n$ and $\Hm_n$ (or the
natural homomorphisms \eqref{eqn:change-of-group-coh}) as the special cases
$\varphi=\identity_G$ (or $f=\identity_M$). These homomorphisms also lift to
the level of derived categories.

% segment-id: o014.li2.chapter6.change-of-group.p012
Equivariant module homomorphisms can also be composed in tandem with group
homomorphisms; the definition is evident.

% segment-id: o014.li2.chapter6.change-of-group.pr001
\begin{proposition}[Functoriality with respect to equivariant homomorphisms]
\label{prop:change-of-group-composite}
	Consider group homomorphisms
	$F\xleftarrow{\psi}G\xleftarrow{\varphi}H$ and equivariant module
	homomorphisms $L\xrightarrow{g}M\xrightarrow{f}N$, where $L$ is an
	$F$-module, $M$ a $G$-module, and $N$ an $H$-module. For every
	$n\in\Z$, we have $\Hm^n(fg)=\Hm^n(f)\Hm^n(g)$.

	If instead the equivariant module homomorphisms point in the reverse
	direction, $L\xleftarrow{g}M\xleftarrow{f}N$, then
	$\Hm_n(gf)=\Hm_n(g)\Hm_n(f)$. Analogous composition equalities also hold
	at the level of derived categories.
\end{proposition}
% segment-id: o014.li2.chapter6.change-of-group.pf003
\begin{proof}
	It suffices to discuss the cohomology case. Fix $n$. By
	Lemma~\ref{prop:change-of-group-compo}, it is enough to prove that the
	diagram
% segment-id: o014.li2.chapter6.change-of-group.g005
	\[\begin{tikzcd}
		\Hm^n(F, L) \arrow[d] & & \\
		\Hm^n(G, \psi^* L) \arrow[r] \arrow[d] & \Hm^n(G, M) \arrow[d] & \\
		\Hm^n(H, \varphi^* \psi^* L) \arrow[r] & \Hm^n(H, \varphi^* M) \arrow[r] & \Hm^n(H, N)
	\end{tikzcd}\]
	commutes. But each square commutes by the functoriality of
	\eqref{eqn:change-of-group-coh}. The derived-category version is entirely
	similar.
\end{proof}

% segment-id: o014.li2.chapter6.change-of-group.ex001
\begin{example}[Conjugation]\label{eg:change-of-group-conj}
	Let $H\lhd G$. Every $t\in G$ induces an automorphism of $H$,
	$a_t:h\mapsto t^{-1}ht$. Now let $M$ be a $G$-module and, to simplify
	notation, continue to write $M$ for $\Res^G_HM$. The homomorphism
	$c_t:M\xrightarrow{x\mapsto tx}M$ is equivariant with respect to the
	oppositely directed group homomorphism $H\xleftarrow{a_t}H$. The
	corresponding homomorphism
% segment-id: o014.li2.chapter6.change-of-group.q011
	\[ \Hm^n(c_t): \Hm^n(H, M) \to \Hm^n(H, M), \]
	decomposes by definition as
	$\Hm^n(H,M)\to\Hm^n(H,a_t^*M)\xrightarrow{c_t}\Hm^n(H,M)$; each part is a
	morphism between cohomological $\delta$-functors.
% segment-id: o014.li2.chapter6.change-of-group.l003
	\begin{itemize}
% segment-id: o014.li2.chapter6.change-of-group.i005
		\item When $n=0$, the first part is simply
		$\identity_M:M^H\to a_t^*(M)^H=M^H$.
% segment-id: o014.li2.chapter6.change-of-group.i006
		\item When $n=0$, the second part is $c_t:M^H\to M^H$, which depends
		only on the coset $tH$.
	\end{itemize}
	If an injective resolution $M\to I$ of the $G$-module is chosen, applying
	$c_t:I^{n,H}\to I^{n,H}$ termwise gives $\Hm^n(c_t)$. When $G=H$, the
	coset $tG=G$, so $\Hm^n(c_t)=\identity$.

	The homology case is similar, but $a_t$ must be replaced by the group
	automorphism $b_t:h\mapsto tht^{-1}$, which points in the same direction
	as the homomorphism $c_t$; the reader is invited to verify this.
\end{example}

% segment-id: o014.li2.chapter6.change-of-group.p013
Finally, we explain how to define the canonical homomorphisms of
Definition~\ref{def:group-equivariance} explicitly using the standard complex
(or standard chain complex) of Proposition~\ref{prop:groupcoh-cochain} (or
Proposition~\ref{prop:grouph-chain}). This makes all the operations concrete.

% segment-id: o014.li2.chapter6.change-of-group.pr002
\begin{proposition}\label{prop:grp-equiv}
	Consider a group homomorphism $\varphi:H\to G$. Let $M$ be a $G$-module,
	$N$ an $H$-module, and $n\in\Z_{\geq0}$. Identify $\Hm^n(G,M)$ (or
	$\Hm_n(G,M)$) with the $\Hm^n$ (or $\Hm_n$) of the standard complex
	$C(G,M)$ (or chain complex $C(G,M)$); do the same after replacing
	$(G,M)$ by $(H,N)$.

% segment-id: o014.li2.chapter6.change-of-group.l004
	\begin{itemize}
% segment-id: o014.li2.chapter6.change-of-group.i007
		\item Let the module homomorphism $f:M\to N$ be
		$\varphi$-equivariant. Then $\Hm^n(f)$ is induced by the morphism of
		complexes
% segment-id: o014.li2.chapter6.change-of-group.g006
		\[\begin{tikzcd}[row sep=tiny]
			C^m(G, M) \arrow[r] & C^m(H, N) \\
			c \arrow[mapsto, r] & {\left[ (h_1, \ldots, h_m) \mapsto f(c(\varphi(h_1), \ldots, \varphi(h_m))) \right]}.
		\end{tikzcd}\]
% segment-id: o014.li2.chapter6.change-of-group.i008
		\item Let the module homomorphism $f:N\to M$ be
		$\varphi$-equivariant. Then $\Hm_n(f)$ is induced by the morphism of
		chain complexes
% segment-id: o014.li2.chapter6.change-of-group.g007
		\[\begin{tikzcd}[row sep=tiny]
			C_m(H, N) \arrow[r] & C_m(G, M) \\
			(h_1 | \cdots | h_m)^\dagger \otimes y \arrow[mapsto, r] & (\varphi(h_1) | \cdots | \varphi(h_m))^\dagger \otimes f(y).
		\end{tikzcd}\]
	\end{itemize}
	In fact, these formulas give the morphism induced by $f$ at the level of
	derived categories.

	The result above also applies, with the same formulas, to the normalized
	standard complex $(\overline{C}^m(G,M))_m$ and normalized chain complex
	$(\overline{C}_m(G,M))_m$ introduced in
	Remark~\ref{rem:nor-cochain}.
\end{proposition}
% segment-id: o014.li2.chapter6.change-of-group.pf004
\begin{proof}
	The problem divides into two parts: first the case $f=\identity_M$, which
	is an explicit description of the canonical homomorphisms
	\eqref{eqn:change-of-group-coh}, and second the case
	$\varphi=\identity$. The latter is straightforward, so we discuss only
	the former.

	Consider the cohomology case. Proposition~\ref{prop:trivial-mod-resolution}
	gives a projective resolution of $G$-modules
	$\mathsf{L}^G\to\Bbbk$ (or a projective resolution of $H$-modules
	$\mathsf{L}^H\to\Bbbk$). Now consider
% segment-id: o014.li2.chapter6.change-of-group.q012
	\[ \Hom^\bullet_G \left( \mathsf{L}^G, \cdot\right) \xrightarrow{\varphi^*} \Hom^\bullet_H\left( \varphi^* (\mathsf{L}^G), \varphi^*(\cdot)\right) \to \Hom^\bullet_H\left(\mathsf{L}^H, \varphi^*(\cdot)\right), \]
	where the second part comes from the morphism
	$\mathsf{L}^H\to\varphi^*(\mathsf{L}^G)$
	(see Lemma~\ref{prop:Shapiro-alpha-beta}). Denote its composite by
	$\Theta$. The asserted morphism of complexes
	$C(G,M)\to C(H,\varphi^*(M))$ corresponds to the action of $\Theta$ on
	the $G$-module $M$.

	We now prove that this formula indeed gives the canonical homomorphism
	defined above. The first method returns to the definition
	\eqref{eqn:change-of-group-coh}, which comes from the universal
	cohomological $\delta$-functor property. It suffices to show that
	$C(G,M)\to C(H,\varphi^*M)$ induces the embedding
	$M^G\to(\varphi^*M)^H$ on $\Hm^0$ and is compatible with long exact
	sequences. Details are left to the reader.

	The second method works in the derived category. The functor
	$\Hom^\bullet_H(\mathsf{L}^H,\cdot)$ preserves acyclic complexes, so it
	induces a functor $\cate{D}^+(H)\to\cate{D}^+(\Bbbk)$ at the level of
	derived categories, with the same notation. This functor can be
	identified with $\RHom_H(\Bbbk,\cdot)$. Thus verifying
	\eqref{eqn:change-of-groups-univ} is equivalent to verifying the following
	equality of vertical composites:
% segment-id: o014.li2.chapter6.change-of-group.g008
	\begin{equation*}
		\begin{tikzcd}[column sep=huge, row sep=large]
			\cate{K}^+(G) \arrow[r, "{\Hom^\bullet_G(\Bbbk, \cdot)}"] \arrow[d] & \cate{K}^+(\Bbbk) \arrow[d] \arrow[Rightarrow, ld] \\
			\cate{D}^+(G) \arrow[r, "{\Hom^\bullet_G(\mathsf{L}^G, \cdot)}", ""' {name=D}] \arrow[r, bend right=60, "" {name=U}, "{\Hom^\bullet_H(\mathsf{L}^H, \varphi^*(\cdot))}"'] & \cate{D}^+(\Bbbk) \arrow[Rightarrow, from=D, to=U, "\Theta"']
		\end{tikzcd} = \begin{tikzcd}[column sep=huge, row sep=large]
			\cate{K}^+(G) \arrow[d] \arrow[r, "{\varphi^*}"] \arrow[rr, bend left, "{\Hom^\bullet_G(\Bbbk, \cdot)}", "{}"' {name=U}] & \cate{K}^+(H) \arrow[d] \arrow[r, "{\Hom^\bullet_H(\Bbbk, \cdot)}"] \arrow[from=U, Rightarrow] \arrow[Rightarrow, ld, "\sim" sloped] & \cate{K}^+(\Bbbk) \arrow[d] \arrow[Rightarrow, ld] \\
			\cate{D}^+(G) \arrow[r, "{\varphi^*}"'] & \cate{D}^+(H) \arrow[r, "{\Hom^\bullet_H(\mathsf{L}^H, \cdot)}"'] & \cate{D}^+(\Bbbk) ,
		\end{tikzcd}
	\end{equation*}
	The symbols \rotatebox[origin=c]{45}{$\Leftarrow$} involving
	$\Hom^\bullet$ come from $\mathsf{L}^G\to\Bbbk$ and
	$\mathsf{L}^H\to\Bbbk$. The complex-level version of the morphism
	$(\cdot)^G\to(\cdot)^H\circ\varphi^*$ is readily identified with
	$\varphi^*:\Hom^\bullet_G(\Bbbk,\cdot)\to
	\Hom^\bullet_H(\varphi^*\Bbbk,\varphi^*(\cdot))
	=\Hom^\bullet_H(\Bbbk,\varphi^*(\cdot))$. The desired equality therefore
	reduces to commutativity of the outer frame of the following diagram of
	complexes:
% segment-id: o014.li2.chapter6.change-of-group.g009
	\[\begin{tikzcd}
		\Hom^\bullet_G(\Bbbk, \cdot) \arrow[r] \arrow[d, "{\varphi^*}"'] & \Hom^\bullet_G(\mathsf{L}^G, \cdot) \arrow[d, "{\varphi^*}"] \arrow[rd, "\Theta"] & \\
		\Hom^\bullet_H(\Bbbk, \varphi^*(\cdot)) \arrow[r] \arrow[rr, bend right=30, "{\text{induced by}\; \mathsf{L}^H \to \Bbbk}"] & \Hom^\bullet_H(\varphi^* \mathsf{L}^G, \varphi^*(\cdot)) \arrow[r] & \Hom^\bullet_H(\mathsf{L}^H, \varphi^*(\cdot)).
	\end{tikzcd}\]
	It remains only to show that the curved lower part commutes, but this is
	exactly the content of Lemma~\ref{prop:Shapiro-alpha-beta}. This proves the
	cohomology case. The homology case is similar.

	Finally, the version for $\overline{C}(G,M)$ is obtained by replacing
	$\mathsf{L}^G$ and $\mathsf{L}^H$ in the proof with
	$\overline{\mathsf{L}}^G$ and $\overline{\mathsf{L}}^H$, respectively;
	everything else is unchanged.
\end{proof}

% segment-id: o014.li2.chapter6.change-of-group.p014
Once the explicit formulas of Proposition~\ref{prop:grp-equiv} are available,
the equalities $\Hm^n(fg)=\Hm^n(f)\Hm^n(g)$ and
$\Hm_n(gf)=\Hm_n(g)\Hm_n(f)$ in
Proposition~\ref{prop:change-of-group-composite}, as well as their
derived-category versions, are immediate at the level of complexes.
